\begin{Theorem}
\label{thm:bessel-parseval-orthonormal-expansion}
设 $\{x_i\}_{i\in I}$ 为 Hilbert 空间 $H$ 中的标准正交集. 那么\begin{enumerate}
\item 对每个 $x\in X$, 有 $\sum_{i\in I}|(x_i,x)|^2\leq\|x\|^2$.
\item 条件 1 中的等号成立的充分必要条件为 $x=\sum_{i\in I}(x_i,x)x_i$.
\end{enumerate}
\end{Theorem}

条件 1 中的不等式称为 Bessel 不等式, 或者在等号成立时称为 Parseval 等式.

\begin{Proof}
对 $I$ 的有限子集 $J$, 利用 $\{x_i\}_{i\in I}$ 是标准正交集这一事实, 得到:\[
\begin{aligned}
0&\leq\left\|x-\sum_{i\in J}(x_i,x)x_i\right\|^2\\
&=\left(x-\sum_{i\in J}(x_i,x)x_i,x-\sum_{l\in J}(x_l,x)x_l\right)\\
&=\|x\|^2-\sum_{i\in J}(x_i,x)(x,x_i)-\sum_{l\in J}(x,x_l)(x_l,x)\\
&\quad+\left(\sum_{i\in J}(x_i,x)x_i,\sum_{l\in J}(x_l,x)x_l\right)\\
&=\|x\|^2-2\sum_{i\in J}|(x_i,x)|^2+\sum_{i\in J}\sum_{l\in J}\overline{(x_i,x)}(x_l,x)(x_i,x_l)\\
&=\|x\|^2-\sum_{i\in J}|(x_i,x)|^2.
\end{aligned}
\]因而 $\sum_{i\in I}|(x_i,x)|^2$ 存在, 且 Bessel 不等式成立. 由 Lemma~\ref{lem:orthogonal-series-summability}, 这意味着 $\sum_{i\in I}(x_i,x)x_i$ 也存在. 此外, 利用上述论证,
\begin{equation}\label{eq:orthonormal-expansion-remainder-norm-identity}
\left\|x-\sum_{i\in I}(x_i,x)x_i\right\|^2=\|x\|^2-\sum_{i\in I}|(x_i,x)|^2,
\end{equation}
该定理的第二个断言立即得出.
\end{Proof}

\begin{Definition}
\label{def:hilbert-basis}
若对每个 $x\in H$ 都有 $x=\sum_{i\in I}(x_i,x)x_i$, 则称 Hilbert 空间 $H$ 中的标准正交集 $\{x_i\}_{i\in I}$ 为一组\textbf{基}.
\end{Definition}

与向量空间基的通常定义不同, 在将 $x$ 表示为 $x_i$ 的线性组合时允许无穷级数. 若可能混淆, 则将 Definition~\ref{def:hilbert-basis} 意义下的基称为 Hilbert 基, 而将有限线性组合意义下的向量空间基称为 Hamel 基. 以下所说的基总是 Hilbert 基.
